\section{The quadratic Poisson reference}
\label{p14:reference:section}

The reconstruction is carried out in the algebra of signed measures.
For the parameters of the model its final reference is a probability.
The following positivity and continuity facts are elementary and do not
depend on the law approximation. The deduction of
Corollary~\ref{p1:cor:tv} from the final theorem is given in
Section~\ref{p14:comp:tv}.

\subsection{Continuity of the reference}

For nonnegative $\mu,\mu'$,
\begin{equation}
 \|\Pi_\mu-\Pi_{\mu'}\|_1\le2|\mu-\mu'|.
 \label{p1:close:poisson-lipschitz}
\end{equation}
For example, if $\mu'=\mu+h$, couple by adding an independent
Poisson variable of mean $h$.  Its probability of being nonzero
is $1-e^{-h}\le h$, proving the inequality.
It follows that
\begin{equation}
 \|Q_{\mu,\alpha}-Q_{\mu',\alpha'}\|_1
 \le2(1+2|\alpha|)|\mu-\mu'|+2|\alpha-\alpha'|.
 \label{p1:close:reference-lipschitz}
\end{equation}
Here we have expanded the difference with coefficient $\alpha$
on $\Pi_\mu-\Pi_{\mu'}$ and used
$\|\Pi_{\mu'}*\Delta^{*2}\|_1\le4$.

\subsection{Positivity and the coefficient of total variation}
\label{p1:close:positivity-section}

The signed representation of the reference is useful in the
proof, but for the parameters of the theorem it is eventually
a probability.  The following exact criterion also specifies
its domain for general parameters.

\begin{lemma}[Positivity of the quadratic Poisson correction]
\label{p1:close:positivity}
Let $\mu>0$, $\alpha\in\mathbb R$, and
$\theta=\mu-\lfloor\mu\rfloor$.  Then $Q_{\mu,\alpha}$ is a
probability measure if and only if
\begin{equation}
 0\le\alpha\le
       \frac{2\mu^2}{\mu+\theta(1-\theta)}.
 \label{p1:close:positivity-range}
\end{equation}
Its masses are
\begin{equation}
 Q_{\mu,\alpha}\{k\}
 =e^{-\mu}\frac{\mu^k}{k!}
   \left[1+\frac{\alpha}{2\mu^2}
                    \bigl((k-\mu)^2-k\bigr)\right],
 \qquad k\ge0.
 \label{p1:close:charlier-masses}
\end{equation}
\end{lemma}

\begin{proof}
Formula \eqref{p1:close:charlier-masses} follows by expressing
the Poisson masses at $k-1$ and $k-2$ relative to the mass
at $k$; negative-index masses are taken to be zero.
The total mass is one because $\Delta^{*2}$ has total mass
zero.

If $\alpha<0$, the bracket is negative for all sufficiently
large $k$, so positivity is impossible.
For $\alpha\ge0$, minimize the quadratic
$f(k)=(k-\mu)^2-k$ over nonnegative integers.
Its successive difference is $f(k+1)-f(k)=2(k-\mu)$.
If $\mu$ is not an integer, the minimum is attained at
$k=\lceil\mu\rceil$; if $\mu$ is an integer, it is attained
at both $k=\mu$ and $k=\mu+1$.  In either case its value is
$-\mu-\theta(1-\theta)$.  Nonnegativity of the bracket
is therefore exactly \eqref{p1:close:positivity-range}.
\end{proof}

In the theorem, $\alpha_N>0$ for large $N$ and $\alpha_N\to0$,
while $\mu_N$ and $\lambda_N$ stay in a compact subinterval
of $(0,\infty)$.  Both corrected references are therefore
probabilities for all sufficiently large $N$.
The norm estimate simultaneously controls all their masses;
it is not only a pointwise assertion at each fixed count.

Define, for $\lambda>0$,
\begin{equation}
 \mathcal C(\lambda)
   =\frac14\|\Pi_\lambda*\Delta^{*2}\|_1
   =\frac14\E\left|
       1-\frac{2P_\lambda}{\lambda}
          +\frac{P_\lambda(P_\lambda-1)}{\lambda^2}
                   \right|,
 \qquad P_\lambda\sim\operatorname{Pois}(\lambda).
 \label{p1:close:charlier-constant}
\end{equation}
The coefficient is explicit even when $\lambda_N$ does not
converge.  The mass at zero and the convolution norm give
\begin{equation}
 \frac{e^{-\lambda}}4\le\mathcal C(\lambda)\le1.
 \label{p1:close:constant-bounds}
\end{equation}
